% ============================================================
% templates/t02-tagged.tex   —— T2 时间标签条目
%
% 适用场景：左边一个短标签、右边一段内容，**标签要对齐成一列**。
%          年份 + 事件、阶段 + 说明、参数名 + 取值、缩写 + 全称。
%          标签之间没有先后关系 —— 有先后关系就用 T9 时间轴。
% 提供    ：\DeckTagged
% 依赖    ：\DeckLine（primitives.tex）、\DeckTagWidth（styles/default.tex）
% 参数    ：#1 标签（≤5 个全角字，超宽会挤到内容上，不报错只难看）
%           #2 内容
% 容量    ：一页最多 \SlideMaxRows 条（与 T1 同）
%           但标签列占了宽度，内容一行只装得下
%           \SlideLineUnitsSafe - 5 个字 —— 所以信息量比 T1 少，
%           别按 T1 的字数往里塞
% 易错    ：① 标签字数不等时**不要用手打空格凑对齐** ——
%             用固定宽度的盒子撑开（\DeckTagged 已经做了这件事）
%           ② 标签每多占一个字，内容就少装一个字
%
% 【填满的样子】
%   \DeckTagged{2023 春}{项目立项，确定以模板库的方式沉淀版式}
%   \DeckTagged{2023 秋}{第一版组件集完成，覆盖七种常见排法}
%   \DeckTagged{2024 春}{引入本地记忆层，新版式不再需要额外授权}
%   \DeckTagged{2024 秋}{吸收机制上线，经验可在 deck 之间复用}
% ============================================================

% --- 带标签的行 ---
% 不用空格去凑对齐（标签字数不等，手打空格永远对不齐）；
% 用固定宽度的盒子把标签撑到统一宽度，内容自然左对齐。
\DeckDefine{\DeckTagged}[2]{%
  \DeckLine{\makebox[\DeckTagWidth][l]{#1}#2}%
}
